% !PRINTED 212
\weberpage{212}
Any rational equation between the roots $x_{\mu,\mu'}$ and $\varkappa^2$ passes by (10) into an identity, even when arbitrary constant numerical coefficients, i.\,e. coefficients independent of $\omega$, occur in it, and one may therefore substitute for $\omega$

\[
\frac{\gamma + \delta\omega}{\alpha + \beta\omega}
\]

i.\,e. one may apply every substitution $T$ to the rational equation between the $x_{\mu,\mu'}$. From this it follows \gesperrt{that the whole group $\mathfrak{B}$ is contained in $\mathfrak{G}$ (Bd.\,I, § 156), and indeed in the group $\mathfrak{G}'$, which arises from the group $\mathfrak{G}$ of the division equation by the adjunction of arbitrary numbers.}

If we extend the domain of rationality of the rational numbers by the adjunction of $n$th roots of unity, then the \gesperrt{Abelian} relations of the preceding section also belong among the rational equations. To these relations, say to

\[
\overset{\nu}{\Sigma} e^{\frac{8\nu\nu'\pi i}{n}} x_{\nu,\nu'} = 0
\]

however, \gesperrt{none} of the substitutions

\[
M = \begin{pmatrix} m, & 0 \\ 0, & 1 \end{pmatrix},
\]

in which $m$ differs from 1, % [source typo] German prints \textquotedbl{}verschwindet\textquotedbl{} (vanishes); \textquotedbl{}verschieden ist\textquotedbl{} is intended.
is applicable; for by this substitution, if $mm' \equiv 1 \pmod{n}$, there passes

\[
\overset{\nu}{\Sigma} e^{\frac{8\nu\nu'\pi i}{n}} x_{\nu,\nu'} \ \text{into} \ \overset{\nu}{\Sigma} e^{\frac{8\nu m'\nu'\pi i}{n}} x_{\nu,\nu'}
\]

which by (12), § 62 is different from zero unless $m'$ and hence also $m \equiv 1 \pmod{n}$. From this follows the theorem:

2. \gesperrt{After adjunction of the $n$th roots of unity (and of arbitrary other constants), $\mathfrak{B}$ is the Galois group of the division equation. $\mathfrak{B}$ is also called the monodromy group of the division equation.}

3. We now prove further that in the original domain of rationality, that is, without adjunction of the $n$th roots of unity or of other irrational numbers, \gesperrt{the group of the division equation is identical with the group $\mathfrak{A}$.}

It is shown in Bd.\,I, § 174 that the primitive $n$th roots of unity $\varrho$ are roots of an irreducible equation

\weq{12}{\Phi(\varrho) = 0}
%CONT>>
